\begin{abstract}
Minimally invasive thoracoscopic and robot-assisted pulmonary surgery reduces operative trauma, but it also removes the direct tactile feedback that surgeons traditionally rely on to localize hidden nodules. Prior tactile studies on breast and other superficial lesions have shown that internal mechanical heterogeneity can be converted into measurable surface responses under controlled contact. Motivated by this transferable physical principle, we investigated whether flexible tactile sensing could be migrated to intraoperative pulmonary nodule localization, despite the additional challenges posed by lung aeration, collapse, deformation, and mismatch between preoperative CT and intraoperative tissue state. We built a tactile prototype system composed of a flexible array sensor, a front-end acquisition unit, and a workstation, and collected an ex vivo porcine-lung dataset with phantom nodules covering seven sizes, six implantation depths, and three repeated trials. We first used mechanism-guided handcrafted features and an XGBoost baseline to progressively verify whether the tactile data supported nodule detection, size inversion, and coarse depth discrimination, and to identify the dominant physical feature families for each task. We then constructed a raw-input hierarchical neural pipeline that learned directly from spatiotemporal tactile tensors in the order of detection, size prediction, and size-aware coarse depth classification. On the independent test set, the raw-input detector achieved an AUC of 0.8383, exceeding the XGBoost baseline of 0.8199. The raw-input size-only model further surpassed XGBoost on size inversion, achieving top-1/top-2/MAE of 0.7177/0.8195/0.1242 versus 0.6701/0.8023/0.1472. For depth, a shared head again collapsed near majority behavior, whereas a size-aware and route-aware raw-input model improved predicted-route balanced accuracy to 0.5240, surpassing the XGBoost depth baseline of 0.5138; a deployment-oriented unified model further improved this value to 0.5337. Explainability analyses, including latent probes, Integrated Gradients, hard-pair analysis, and phase occlusion, indicated that the network encoded part of the physically meaningful structure related to size and depth, although its temporal strategy still relied heavily on the peak neighborhood. Overall, the study supports a progressive conclusion: tactile spatiotemporal sensing can first support reliable intraoperative pulmonary nodule detection, then size inversion, and under explicit hierarchical organization, cautious coarse depth discrimination. The main contribution lies in establishing a validated, interpretable, and engineering-oriented technical route rather than claiming a finalized clinical system.
\end{abstract}
